\subsection{PRODUCT OF CONNECTED SPACES}

PROPOSITION 8. Every product of connected spaces is connected. Conversely, if a product of non-empty spaces is connected, then each of the factors is connected.

Let \(\mathbf{X} = \prod_{i \in I} \mathbf{X}_i\) be a product of topological spaces. If the \(\mathbf{X}_i\) are non-empty, we have \(\mathbf{X}_i = \text{pr}_i\mathbf{X}\) for each \(i \in I\) ; hence if \(\mathbf{X}\) is connected so are the \(\mathbf{X}_i\) (no. 2, Proposition 4). Conversely, suppose that each of the \(\mathbf{X}_i\) is connected and \(\mathbf{X}\) is not. By Proposition 5 of no. 2, there exists a continuous surjective mapping \(f: \mathbf{X} \to \mathbf{X}'\) , where \(\mathbf{X}'\) is a discrete space which contains more than one point. Let \(a = (a_i)\) be any point of \(\mathbf{X}\) , and \(x\) any index; the partial mapping \(f_x: \mathbf{X}_x \to \mathbf{X}'\) , defined by \(f_x(x) = f((y_i))\) where \(y_x = x\) and \(y_i = a_i\) if \(i \neq x\) , is continuous on \(\mathbf{X}_x\) ; since \(\mathbf{X}_x\) is connected, \(f_x\) must be constant on \(\mathbf{X}_x\) . It follows immediately by induction that \(f(x) = f(a)\) for all points \(x = (x_i)\) such that \(x_i = a_i\) for all but a finite number of indices \(i \in I\) . But these points \(x\) form a dense subset of \(\mathbf{X}\) ( \(\S 4\) , no. 3, Proposition 8). Hence \(f\) is continuous on \(\mathbf{X}\) and constant on a dense subset of \(\mathbf{X}\) , and therefore constant on \(\mathbf{X}\) ( \(\S 8\) , no. 1, Proposition 2, Corollary 1). But this contradicts the definition of \(f\) .

